C'est le chapitre le plus rétif du collège, et il faut dire pourquoi avant de commencer.

Au collège, l'aire d'un disque \emph{est} $\pi r^2$ : c'est une formule qu'on donne, qu'on
applique, et qu'aucun élève ne démontre. En mathématiques, c'est un théorème de la théorie
de la mesure — l'aire est une mesure, le disque une partie du plan, et l'égalité demande une
intégration. Formaliser honnêtement la formule d'aire du disque, ou le volume de la boule,
suppose donc un appareillage sans commune mesure avec le niveau de l'énoncé.

Le chapitre démontre ce qui relève du calcul — les conversions, la proportionnalité des
grandeurs, le périmètre — et laisse les formules d'aire et de volume à leur statut de
définitions données. Les trois énoncés concernés portent la marque \og{} en cours \fg{} dans
l'index, avec cette raison.

Ce n'est pas un échec : c'est le résultat. Une bonne partie de ce qu'on appelle \og{} savoir
sa formule \fg{} au collège n'est pas un théorème que l'élève pourrait démontrer, mais une
convention qu'il doit accepter — et la formalisation le rend visible.
